% S06-M2 (package abhijit/s06-math): setting, Hecke algebras, old/new decomposition, Tamagawa ideal.
% Labels: ss:s06-setting, hyp:s06-H, rem:s06-CR, ss:s06-rings, def:s06-rings, lem:s06-decomp, def:s06-ctam, def:s06-tam,
% lem:s06-tam.
\section{Setting: Hida families at a prime \texorpdfstring{$q\,\|\,N$}{q || N}}

\subsection{Setting and standing hypotheses}\label{ss:s06-setting}
Here $D=q$, so $B=B_q$ is ramified at $\{q,\infty\}$. $p\ge5$, $q\ne p$, $(M,pq)=1$, $N:=qM$. $\Lam=\cOE[[1+p\Zp]]$, with
uniformiser $\varpi$ of $\cOE$ and fraction field $K=\mathrm{Frac}\,\Lam$. On the $\GL_2$ side, $\Lam$ acts through the
\emph{classical} diamond operators, and $P_k=\ker([u]\mapsto u^{k-2})$ is the weight-$k$ point (seed \S2.6). On the definite side,
stage 04 lets $\Lam$ act by \emph{adelic} diamonds. These are inverse to the classical ones (stage 01). In those conventions the
weight-$k$ point is $P^{\rm ad}_k=\ker([u]\mapsto u^{2-k})$ (\texttt{if:s04-Mord}, stage 04 D1). The $\Lam$-adic JL
$\T'\cong\T^{\ord}_B$ is $\Lam$-linear only if the two actions are matched by $[u]\mapsto[u]^{-1}$. We always use the $\GL_2$
normalisation and transport it. We write $S_q\in\Lam^\times$ for the element with $P_k(S_q)=q^{k-2}$ (trivial character), so that
$qS_q$ interpolates $q^{k-1}$.

\begin{hypothesis}[standing hypotheses; seed \S9.1, corrected]\label{hyp:s06-H}
\leavevmode
\begin{itemize}
\item[(H0)] $p\ge5$, $q\ne p$, $(M,pq)=1$, $N=qM$, $p\nmid\varphi(N)$. In particular $q\not\equiv1\bmod p$.
\item[(H0$^-$)] (\emph{new}; needed for the fixed-weight input Cor.~\ref{cor:s06-fixed} and for Prop.~\ref{prop:s06-betti};
  Thm~\ref{thm:s06-defect} needs (H0$^-$) \emph{or} (Lev); the inclusion of Prop.~\ref{prop:s06-galois} uses neither)
  $q\not\equiv-1\bmod p$, or $\bar\rho_\m$ is ramified at $q$. Together with (H0): $p\nmid q^2-1$, or $\bar\rho_\m$ is ramified at $q$.
\item[(Lev)] (\emph{new}; alternative to (H0$^-$) in Thm~\ref{thm:s06-defect}) At the point $P$ of (P), of weight $k=k_P$, with
  $f=F_P$ (ordinary): either (F0) $M=1$; or (F1) $M>1$ is squarefree, $\bar\rho$ is unramified at every $\ell\mid M$ (so $N(\bar\rho)=1$),
  and $\ell\not\equiv-1\bmod p$ for every $\ell\mid M$; or (F1$'$) as (F1), but with some $\ell\mid M$, $\ell\equiv-1\bmod p$. These are the
  cases (F0), (F1), (F1$'$) of Prop.~\ref{prop:s06-levM}, in which its hypothesis (Lev$'_M$) (with (L3)) holds automatically. Status:
  Prop.~\ref{prop:s06-levM} is \stpmc{} in all three cases, for (F1$'$) in our sign-fixed reading R4 of \cite[Thm~3.14]{KO} (ours, not
  \cite{KO}'s text). Not included: $N(\bar\rho)>1$ and non-squarefree $M_0$, which are \stopen{} (Remark~\ref{rem:s06-levM-status}); (L3) is the condition under
  which the method may extend.

\item[(H1)] $\T$ is Hida's full ordinary cuspidal Hecke algebra of tame level $\Gamma_0(N)$ over $\Lam$, in a fixed weight component
  $k_0\bmod(p-1)$. It is generated by all $T_\ell$ ($\ell\nmid Np$), $U_\ell$ ($\ell\mid Np$, in particular $U_q$ and $U_p$), and it
  is localised at a non-Eisenstein maximal ideal $\m$. $F$ is a primitive ordinary $q$-new branch, with $\lambda_F:\T\to I_F$, and
  $I_F$ is a normal local domain finite over $\Lam$. Then $I_F$ is $\Lam$-free: it is normal of dimension $2$, hence
  Cohen--Macaulay (Serre's criterion), hence free over the regular local ring $\Lam$ (Auslander--Buchsbaum). So the seed's
  ``hence $\Lam$-free'' is correct; it is used only through flatness of $R=\T\otimes_\Lam I_F$ over $I_F$.
\item[(H2)] $\bar\rho=\bar\rho_\m$ is absolutely irreducible on $G_{\Q(\sqrt{p^*})}$ and $p$-distinguished.
\item[(Im)] $p\mid\#\bar\rho(G_\Q)$. \emph{None of the cited theorems below uses (Im).} \cite[Thm~1.1]{KO} needs only (H2) (its
  condition (2) $=$ Assumption 2.1(TW)), and \cite[Rem.~2.10]{KO} says that the \cite{DFG} and \cite{Dim09} image conditions are interchangeable
  over $\Q$. We keep (Im) as in the seed, flagged as probably superfluous (Risk R4).
\item[(P)] Either (P1) $k_0\not\equiv2\bmod(p-1)$, and then $F$ has an arithmetic point $P$ of weight $k_P\in[3,p-1]$ with trivial
  character (a \emph{Kim--Ota point}). Or (P2) $k_0\equiv2$, $F$ has a weight-two point $P$ which is the ordinary
  $p$-stabilisation of a newform $f$ of level $N$, \emph{and} (\emph{new}) the weight-two ordinary Hecke algebra
  $\T^{\ord}_2(\Gamma_0(Np))_\m$ has no $p$-new eigenforms. The extra condition makes $\T/P\T\cong\T_2(\Gamma_0(N))_\m$,
  which is the algebra of \cite{KO}. Without it the control step (3) of the proof compares with an algebra that \cite{KO} does not treat.
\end{itemize}
\sthyp
\end{hypothesis}

\begin{remark}[relation to Hsieh's residual condition]\label{rem:s06-CR}
Hsieh's (CR,$\Sigma$) is \lit{hsieh\_2017\_hida-families-and-p-adic-triple-product-l-functions\_arxiv}{\S1.4, Hypothesis
(CR,$\Sigma$), p.~4}: $\bar\rho_F$ is absolutely irreducible and $p$-distinguished, and it is ramified at every $\ell\in\Sigma$ with
$\ell\equiv1\bmod p$. The congruence is $+1$ only, not $\pm1$ (stage 04 point (4)). With $\Sigma=\Sigma^-=\{q\}$, (H2)$+$(H0) imply
(CR,$\{q\}$), because absolute irreducibility on $G_{\Q(\sqrt{p^*})}$ implies it on $G_\Q$ and (H0) gives $q\not\equiv1$. Hence the
\emph{combinatorial half needs only (H0), (H1), (H2)}. The congruence $q\equiv-1$ enters only through the fixed-weight Galois input
(Remark~\ref{rem:s06-minus1}). \stproved{} (the implication; the hypothesis itself is \stcited.)
\end{remark}

\subsection{Hecke algebras, congruence ideals, defect}\label{ss:s06-rings}
\begin{definition}[seed \S9.1]\label{def:s06-rings}
$\T'$ and $\T''$ are the images of $\T$ in the endomorphisms of the $q$-new and of the $q$-old ordinary $\Lam$-adic cusp forms
(localised at $\m$; see the proof of Lemma~\ref{lem:s06-decomp}(a) for the two spaces). Put $\fa:=\ker(\T\to\T')$ and
$\fb:=\ker(\T\to\T'')$. Put $R:=\T\otimes_\Lam I_F$, $R':=\T'\otimes_\Lam I_F$, $R'':=\T''\otimes_\Lam I_F$, and keep the names
$\fa,\fb$ for $\fa R,\fb R$. $\lambda:R\to I_F$ is the $I_F$-linear extension of $\lambda_F$. Since $F$ is $q$-new, $\lambda$ factors
as $\lambda=\lambda'\circ(R\to R')$. For a finite $A$-algebra $S$ and an $A$-algebra map $\mu:S\to A$ to a domain,
$\eta(S,\mu):=\mu(\Ann_S\ker\mu)$ (seed \S2.6; Definition~\ref{def:s06-gor}). Put
\[\eta_F:=\eta(R,\lambda),\qquad \eta^B_F:=\eta(R',\lambda'),\qquad (h_F):=\eta_F\,(\eta^B_F)^{-1}\ \text{ when both are principal.}\]
\status{definition}
\end{definition}
The seed writes $h_F=\eta_F/\eta^B_F$ as an element. This makes sense only when both ideals are principal, which needs (Gor$_R$) and
(Gor$_{R'}$) (Lemma~\ref{lem:s06-S}), and $\eta^B_F\ne0$. Under them, $h_F\in I_F$ is well defined up to $I_F^\times$
(Thm~\ref{thm:s06-E}).

\begin{lemma}[old/new decomposition]\label{lem:s06-decomp}
Assume (H1). (a) $\T\otimes_\Lam K=(\T'\otimes K)\times(\T''\otimes K)$ with $K=\mathrm{Frac}\,\Lam$, $\fa\cap\fb=0$, and
$\T/\fa$, $\T/\fb$ are $\Lam$-torsion-free. (b) Let $\T(M)_\m$ be Hida's ordinary algebra of tame level $M$ localised at the
pull-back of $\m$, and $S_q\in\Lam^\times$ the element with $P_k(S_q)=q^{k-2}$ (trivial character). If $\bar\rho_\m$ is ramified at
$q$, then $\T''=0$. If $\bar\rho_\m$ is unramified at $q$ and $\T''\ne0$ (so that the pull-back $\m_M$ of $\m$ to $\T(M)$ is defined; if $\T''=0$ put $\T(M)_\m:=0$), $X^2-T_qX+qS_q$ has two roots in $\T(M)_\m$ that are distinct modulo the maximal ideal: they are
$\equiv U_q$ and $\equiv qU_q$ mod $\m$, and $q\not\equiv1$ by (H0). Let $\alpha_q$ be the root congruent to $U_q$. Then
$\T''\cong\T(M)_\m$ via $U_q\mapsto\alpha_q$. In particular $\T''$ is $\Lam$-free, and $T_q=U_q+qS_qU_q^{-1}$ lies in the image
of $\T$. The same holds at fixed weight $k$: $\T''_k\cong\T_k(M)_\m$, $U_q\mapsto\alpha_q$.
\provisional{F2: the old quotient is the localisation $\T(M)_\m$, not $\T(M)_\m[X]/(X^2-T_qX+qS_q)$ (which has twice the rank);
before: ``$\T''\cong\T(M)_\m[X]/(X^2-T_qX+qS_q)$, $X\mapsto U_q$; in particular $\T''$ is $\Lam$-free''.} (c) The same statements
hold after $\otimes_\Lam I_F$ for $R,R',R''$; in particular $(R,\fa,\fb)$ satisfies the hypotheses of Lemma~\ref{lem:s06-D} over
$A=I_F$, and $R''=R/\fb$ is $I_F$-free. \stpmc{} (proof below; cited inputs: Hida freeness and control, local--global
compatibility at $q$.)
\end{lemma}

\noindent\emph{Cited inputs.}
(I1) \cite[Thm~3.1, Cor.~3.2]{Hi86b}. $\T$ and $\T(M)$ are free of finite rank over $\Lam$ ($p\ge5$): \cite[Thm~3.1, p.~250]{Hi86b}
\lit{hida\_1986\_iwasawa-modules-attached-to-congruences-of-cusp-forms}{Thm 3.1, p.~250} \stcited. Hida states it for the
ordinary parts $h^0(N;\cO_K)$ of the universal Hecke algebras of level $\Gamma_1(Np)$; our algebras are the trivial-tame-character,
weight-component and $\m$-local direct factors (a direct factor of a free module over the local ring $\Lam$ is free), and ordinary
$\Lam$-adic cusp forms are the $\Lam$-dual of $\T$ by the duality of \cite[\S2]{Hi86b} (so they form a $\Lam$-free module $\mathbf S$ on
which $\T$ acts faithfully). Control: \cite[Cor.~3.2, (3.3b), p.~250]{Hi86b} \stcited.
(I2) \cite[Thm~3.1]{DDT97}, \cite[Th.~(A)]{Car86}. For a newform $f$ of weight $k\ge2$, level $N_f$ and trivial character, and a prime $\ell\ne p$: if $\ell\nmid N_f$ then $\rho_f$
is unramified at $\ell$; if $\ell\,\|\,N_f$ then $\rho_f|_{G_\ell}\cong\big(\begin{smallmatrix}\chi\epsilon&*\\0&\chi\end{smallmatrix}\big)$ with
$\chi$ unramified, $\chi(\Frob_\ell)=a_\ell(f)$, $\epsilon$ the cyclotomic character and $*$ non-trivial on inertia (the conductor
exponent at $\ell$ is $1$); since $\det\rho_f=\epsilon^{k-1}$, this gives $a_\ell(f)^2=\ell^{k-2}$:
\lit{darmon-diamond-taylor\_1997\_fermats-last-theorem\_preprint}{Thm 3.1(d),(e), p.~85--86 (weight 2)} and in general
\lit{carayol\_1986\_sur-les-representations-l-adiques-associees-aux-formes-modulaires-de-hilbert}{Th\'eor\`eme (A), \S0.7,
p.~410} \stcited.

\begin{proof}
(a) \emph{The two spaces.} Let $\mathbf S$ be the $\Lam$-module of ordinary $\Lam$-adic cusp forms of tame level $N$ in our component,
localised at $\m$ (I1). The commutative $K$-algebra $\T_K:=\T\otimes_\Lam K$ is finite-dimensional, hence a finite product of local
Artinian algebras $\T_{K,\mathfrak n}$; correspondingly $\mathbf S_K=\bigoplus_{\mathfrak n}\mathbf S_{K,\mathfrak n}$ with $\T_{K,\mathfrak n}$ acting
faithfully on $\mathbf S_{K,\mathfrak n}$ (because $\T$ acts faithfully on $\mathbf S$). The maximal ideals $\mathfrak n$ correspond to
$\mathrm{Gal}(\bar K/K)$-orbits of $\Lam$-algebra maps $\phi:\T\to\bar K$, i.e.\ to normalised ordinary $\Lam$-adic eigenforms, and each
$\phi$ has a Galois representation $\rho_\phi$ over a finite extension of $K$, characterised by $\tr\rho_\phi(\Frob_\ell)=\phi(T_\ell)$
for $\ell\nmid Np$ (as in Def.~\ref{def:s06-tam}). Call $\mathfrak n$ \emph{$q$-new} if $\rho_\phi$ is ramified at $q$, and \emph{$q$-old}
otherwise. (The existence of $\rho_\phi$ is \cite[Thm~2.1]{Hi86a}; \stcited{}.) At a Zariski-dense set of arithmetic points where $\phi$ specialises to (the $p$-stabilisation of) a classical newform $f$ (all that is used), (I2) shows: $q$-old
$\iff$ $q\nmid N_f$ $\iff$ $f$ comes from level $M$. Put $\mathbf S^{\rm new}:=\mathbf S\cap\bigoplus_{\text{new}}\mathbf S_{K,\mathfrak n}$ and
$\mathbf S^{\rm old}:=\mathbf S\cap\bigoplus_{\text{old}}\mathbf S_{K,\mathfrak n}$. These are $\T$-stable, $\Lam$-saturated, and
$\mathbf S_K=\mathbf S^{\rm new}_K\oplus\mathbf S^{\rm old}_K$. They are the spaces of Definition~\ref{def:s06-rings}: $\mathbf S^{\rm old}_K$ is
spanned by $G(z)$, $G(qz)$ for $G$ of tame level $M$ (proof of (b)).

\emph{The product decomposition.} The image of $\T_K$ in $\End_K(\mathbf S^{\rm new}_K)$ is $\prod_{\text{new}}\T_{K,\mathfrak n}$, and in
$\End_K(\mathbf S^{\rm old}_K)$ it is $\prod_{\text{old}}\T_{K,\mathfrak n}$. Hence $\T'\otimes K=\prod_{\rm new}\T_{K,\mathfrak n}$,
$\T''\otimes K=\prod_{\rm old}\T_{K,\mathfrak n}$ (localisation is exact, so the image of $\T$ tensored with $K$ is the image of
$\T_K$), and $\T_K=\T'_K\times\T''_K$, because every $\mathfrak n$ is either new or old and not both. (The content of ``not both'' is the
local statement (I2): the Galois representation is either ramified or unramified at $q$. This replaces the strong-multiplicity-one
argument of the tex.)

\emph{$\fa\cap\fb=0$.} An element of $\fa\cap\fb$ kills $\mathbf S^{\rm new}+\mathbf S^{\rm old}$, which has full rank in $\mathbf S$; since
$\mathbf S$ is $\Lam$-free it kills $\mathbf S$, and it is $0$ by faithfulness.

\emph{Torsion-freeness.} $\T/\fa\cong\T'\subseteq\End_\Lam(\mathbf S^{\rm new})$, and $\mathbf S^{\rm new}\subseteq\mathbf S$ is torsion-free, hence
so is $\End_\Lam(\mathbf S^{\rm new})$. The same for $\T/\fb$.

(b) \emph{The ramified case.} If $\T''\ne0$, some $q$-old $\phi$ occurs at $\m$; then $\bar\rho_\m\cong\bar\rho_\phi^{\rm ss}$ (Brauer--Nesbitt
and Chebotarev, using (H2)) is unramified at $q$. So $\bar\rho_\m$ ramified at $q$ forces $\T''=0$.

\emph{The old space.} Let $\mathbf S(M)$ be the ordinary $\Lam$-adic cusp forms of tame level $M$ (same component). The degeneracy
maps $G\mapsto G(z)$, $G\mapsto G(qz)$ commute with $e$ (as $q\ne p$) and give a map $\iota:\mathbf S(M)^2\to\mathbf S$ (before
localisation) whose image spans the $q$-old part over $K$: every $q$-old $\phi$ is a $\Lam$-adic eigenform whose primitive form has tame
level prime to $q$, i.e.\ dividing $M$, and conversely. On $\mathbf S(M)^2$ the operators of level $N$ act as follows: $T_\ell$
($\ell\nmid Np$), $U_\ell$ ($\ell\mid Mp$) diagonally by the level-$M$ operators, and $U_q$ by the matrix
$\big(\begin{smallmatrix}T_q&1\\-qS_q&0\end{smallmatrix}\big)$ in the basis $(G(z),G(qz))$, since $U_q(G)=T_qG-qS_q\,G(qz)$ and
$U_q(G(qz))=G$ (the classical identities $a_{qn}(G)=a_q(G)a_n(G)-q^{k-1}a_{n/q}(G)$ at weight $k$, interpolated). So the action of
the level-$N$ operators on $\iota(\mathbf S(M)^2)_K$ factors through $\T(M)[X]/(X^2-T_qX+qS_q)$, $X\mapsto U_q$, acting on
$\mathbf S(M)\otimes_{\T(M)}\T(M)[X]/(f)$, $f:=X^2-T_qX+qS_q$.

\emph{The roots mod $\m$.} Write $u:=U_q\bmod\m\in k_\m^\times$ ($U_q$ is a unit: on $q$-new forms $U_q^2=S_q$, a unit).
By (H1) $F$ is a $q$-new branch in the $\m$-component, so $\m$ contains $U_q^2-S_q$ (for $q$-new eigensystems with trivial character
and $q\,\|\,N$, $a_q^2=q^{k-2}$ at every arithmetic point, i.e.\ $U_q^2=S_q$ on $\T'$), and $\bar S_q=u^2$. If $\T''\ne0$, the pull-back
$\m_M$ of $\m$ to $\T(M)$ (the maximal ideal with the same residual eigenvalues of $T_\ell$, $\ell\nmid Np$, $U_\ell$, $\ell\mid Mp$, and
$T_q\mapsto u+q\bar S_qu^{-1}$) is defined, and mod $\m_M$ the polynomial $f$ has the root $u$ (as $U_q$ acts on the old part by a root of
$f$). The other root is $q\bar S_q/u=qu^2/u=qu$, and $qu\ne u$ because $q\not\equiv1\bmod p$ (H0). This uses $\T'\ne0$, i.e.\ the
existence of $F$; without a $q$-new form at $\m$ the two roots can coincide mod $\m$.

\emph{Hensel.} $\T(M)_\m:=\T(M)_{\m_M}$ is complete local (finite over $\Lam$, I1). The discriminant $d=T_q^2-4qS_q$ reduces to
$(u-qu)^2\ne0$, a unit, and $p\ne2$; Newton's iteration $x_{n+1}=\tfrac12(x_n+d/x_n)$ from a lift of $u-qu$ converges $\m$-adically to a
square root of $d$. So $f=(X-\alpha_q)(X-\beta_q)$ with $\alpha_q\equiv u$, $\beta_q\equiv qu$, $\alpha_q-\beta_q\in\T(M)_\m^\times$, and by the
Chinese remainder theorem $\T(M)_\m[X]/(f)\cong\T(M)_\m\times\T(M)_\m$, $X\mapsto(\alpha_q,\beta_q)$.

\emph{Localising at $\m$.} The $\m$-part of the old space is the part where $U_q\equiv u$, i.e.\ the first factor:
$\mathbf S^{\rm old}_K=\mathbf S(M)_{\m_M,K}$ with $U_q$ acting as $\alpha_q$ (the second factor has $U_q\equiv qu\not\equiv u$ and lies in
another maximal ideal of $\T$). Hence $\T''$ is the closed $\Lam$-subalgebra of $\End(\mathbf S(M)_{\m_M})$ generated by the images of
$T_\ell$ ($\ell\nmid Np$), $U_\ell$ ($\ell\mid Mp$) and $\alpha_q$; it lies in $\T(M)_\m$ (which acts faithfully on $\mathbf S(M)_{\m_M}$). It
contains $T_q=\alpha_q+\beta_q=\alpha_q+qS_q\alpha_q^{-1}$, so it contains the images of all generators of $\T(M)$, hence all of
$\T(M)_\m$ (a direct factor of $\T(M)$, I1). So $\T''\cong\T(M)_\m$ with $U_q\mapsto\alpha_q$, it is $\Lam$-free by (I1), and $T_q$ is the
image of $U_q+qS_qU_q^{-1}\in\T$.

\emph{Fixed weight.} The same argument with $\cO$ in place of $\Lam$, the classical spaces of weight $k$ in place of $\mathbf S$ and
$q^{k-2}$ in place of $S_q$ gives $\T''_k\cong\T_k(M)_\m$, $U_q\mapsto\alpha_q$ (the root of $X^2-T_qX+q^{k-1}$ congruent to $U_q$).

(c) $I_F$ is $\Lam$-free (H1), so $-\otimes_\Lam I_F$ is exact. Hence $R$ is finite free over $I_F$; $\fa R=\fa\otimes I_F$ and
$\fb R=\fb\otimes I_F$; $\fa R\cap\fb R=(\fa\cap\fb)\otimes I_F=0$; $R/\fa R=\T'\otimes I_F$ embeds into $\T'_K\otimes_\Lam I_F=\T'\otimes_\Lam
K_F$ ($K\otimes_\Lam I_F=K_F:=\mathrm{Frac}\,I_F$ as $I_F$ is finite over $\Lam$), a $K_F$-vector space, so it is $I_F$-torsion-free; the
same for $R/\fb R$, which is moreover $I_F$-free since $\T''$ is $\Lam$-free. Finally an element $0\ne c\in\Lam$ with $c\in\fa+\fb$
(which exists by (a) and Lemma~\ref{lem:s06-split}) stays non-zero in $I_F\supseteq\Lam$.
\end{proof}
\noindent\emph{Remarks on the sources.} $\Lam$-freeness of $\T(M)_\m$ is \cite[Thm~3.1]{Hi86b}
\lit{hida\_1986\_iwasawa-modules-attached-to-congruences-of-cusp-forms}{Thm 3.1, p.~250} ($p\ge5$). The page reference is
re-checked against the OCR: Thm~3.1, Cor.~3.2 (with (3.3a), (3.3b)), the definition of old forms (3.4a--d) and Cor.~3.3 are all on
journal p.~250 (PDF p.~21; journal page $=$ PDF page $+229$). \cite[Cor.~3.3, p.~250]{Hi86b} gives $\Lam$-freeness of the old
\emph{ideal}, not of $\T''$. Hida's old ideal is the annihilator of the forms old at \emph{any} divisor $t>1$ of the tame level
\cite[(3.4a--d), p.~250]{Hi86b}, so it is contained in our $\fb$ ($q$-old only) and equals it only when no form in the $\m$-part is old at a
prime of $M$ (e.g.\ $M=1$). \provisional{F2: Hensel/localisation argument replaces the Chebotarev argument for ``$T_q$ in the
image''; before: the subtle point was surjectivity onto $\T(M)_\m[X]/(\cdots)$, proved by topological generation of $\T(M)_\m$ by the
$T_\ell$, $\ell\ne q$ (Chebotarev).} \provisional{S06-M2: the distinctness of the two roots mod $\m$ (``$\equiv U_q$ and $\equiv qU_q$'') uses
$U_q^2\equiv S_q\bmod\m$, which holds because the $q$-new branch $F$ lies in the $\m$-component (H1); the tex states it without this
justification; statement unchanged.}

\begin{definition}[combinatorial Tamagawa ideal; seed Def.~ctam]\label{def:s06-ctam}
$\cc:=\mathrm{im}(\fb\to R')=(\fa+\fb)/\fa$, so $R'/\cc\cong R/(\fa+\fb)$ is the $\Lam$-adic old/new congruence module at $q$
(after $\otimes I_F$). Put $\cc_q(F):=\lambda'(\cc)=\lambda(\fb)\subseteq I_F$. At an arithmetic point $P$ with $q$-new
specialisation $f$, the \emph{fixed-weight} version is $\cc_q(f):=\lambda_f(\fb_k)$, where $\fb_k=\ker(\T_k(N)_\m\to\T''_k)$ is taken
in the weight-$k$ Hecke algebra. \status{definition}
\end{definition}
\noindent This is Diamond's $\Lam$-adic old/new congruence module \cite{Dia89}. Diamond computes its specialisation along the \emph{old}
branch: \cite[Thm~6B, p.~80]{Dia89}
\lit{diamond\_1989\_congruence-modules-associated-to-lambda-adic-forms}{Thm 6A--6C, p.~80} \stcited{} gives
$\mathrm{div}(\mathfrak b)=-\mathrm{div}(a_q-\xi(q+1))$ for the sign-$\xi$ part. Here we use the \emph{new} branch, where the same
module is read through $\lambda'$ instead.

\begin{definition}[Tamagawa ideal; seed \S9.1]\label{def:s06-tam}
Let $H$ be Hida's ordinary $\Lam$-adic étale cohomology of tame level $N$ localised at $\m$. It is free of rank two over $\T$
under (H2) (input (1) of Prop.~\ref{prop:s06-galois}). Put $\mathcal T_F:=H\otimes_{\T,\lambda_F}I_F$, with
$\rho_F:G_\Q\to\GL_2(I_F)$. Define $\tq_q(F)\subseteq I_F$ as the ideal generated by the entries of $\rho_F(\sigma)-1$,
$\sigma\in I_q$. (We write $\mathcal T_F$ rather than the seed's $\cT_F$, since $\cT$ is the tree.) The ideal does not depend on
the basis of $\mathcal T_F$. \status{definition}
\end{definition}

\begin{lemma}\label{lem:s06-tam}
(a) $\rho_F|_{I_q}$ is unipotent and factors through the $\Zp(1)$-quotient of tame inertia. So $\tq_q(F)$ is generated by the
entries of $\rho_F(\tau)-1$ for one generator $\tau$. (b) At an arithmetic point $P$ with $q$-new $f=F_P$,
$P(\tq_q(F))\,\widetilde\cO_P=(\varpi_P^{t_f(q)})$, where $\widetilde\cO_P$ is the normalisation of $I_F/P$ in its fraction field and $P(\tq_q(F))\widetilde\cO_P$ is the ideal generated by the image of $\tq_q(F)$ under $I_F\to I_F/P\subseteq\widetilde\cO_P$ (see the remark after the proof), and $t_f(q)$ is the Tamagawa exponent of
\cite[Def.~2.12]{KO} \lit{kim-ota\_2019\_quantitative-variation-of-congruence-ideals-and-integral-periods-of-modular-forms\_arxiv}{Def.~2.12,
p.~6} (= \cite[Def.~3.3]{PW} \lit{pollack-weston\_2006\_anticyclotomic-mu-invariants-of-modular-forms\_arxiv}{Def.~3.3, p.~9}),
computed with coefficient ring $\widetilde\cO_P$. (c) $t_f(q)<\infty$ because $f$ is $q$-new. \stpmc{} (proof below; inputs (I2) and
uniqueness of the lattice under (H2).)
\end{lemma}
\begin{proof}
(a) \emph{Unipotence.} Let $\sigma\in I_q$. At every arithmetic point $P$ with $q$-new specialisation $f=F_P$, $\mathcal T_F/P$ is a
$G_\Q$-stable lattice in $V_f$, and $\rho_f(\sigma)$ is unipotent by (I2) ($q\,\|\,N_f$). So $(\rho_F(\sigma)-1)^2\in P\,M_2(I_F)$ for a
Zariski-dense set of height-one primes $P$, whose intersection is $0$ in the Noetherian domain $I_F$ (a non-zero element lies in only
finitely many minimal primes over it). Hence $(\rho_F(\sigma)-1)^2=0$.
\emph{Wild inertia.} For $\sigma\in P_q$ (pro-$q$), $\sigma^{q^n}\to1$, so $\rho_F(\sigma)^{q^n}=1+q^n(\rho_F(\sigma)-1)\to1$; as $q\in I_F^\times$,
this forces $\rho_F(\sigma)=1$. \emph{Tame inertia.} $I_q/P_q\cong\prod_{\ell\ne q}\Z_\ell(1)$ is procyclic, so $\rho_F(I_q)$ is the closure of
$\{\rho_F(\tau)^n\}=\{1+nN\}$, $N:=\rho_F(\tau)-1$, which is $\{1+aN:a\in\Zp\}$, a pro-$p$ group (the additive group of the
$\m$-adically complete $I_F$ is pro-$p$). So $\rho_F|_{I_q}$ factors through the $\Zp(1)$-quotient, and the entries of
$\rho_F(\tau^a)-1=aN$ lie in the ideal generated by the entries of $N$.

(b) $P(\tq_q(F))\widetilde\cO_P$ is generated by the entries of $(\rho_F(\tau)-1)\bmod P$, i.e.\ of $\rho_f(\tau)-1$ on the lattice
$T:=\mathcal T_F\otimes_{I_F,P}\widetilde\cO_P$ (the extension to $\widetilde\cO_P$ of the $I_F/P$-lattice $\mathcal T_F/P$; same ideal of entries). By (H2) $\bar\rho_f$ is irreducible, so $T$ is the unique $G_\Q$-stable lattice up to scaling, i.e.\ the lattice
$T_f$ of \cite[\S2.2]{KO}. By \cite[Def.~2.12]{KO}, $t_f(q)$ is the largest $t$ such that $T_f/\varpi_P^tT_f$ is unramified at $q$, i.e.\ such that
$\rho_f(\tau)\equiv1\bmod\varpi_P^t$ ((a) at $P$), i.e.\ such that all entries of $\rho_f(\tau)-1$ lie in $(\varpi_P^t)$. As $\widetilde\cO_P$ is a DVR (remark below),
the ideal generated by the entries is $(\varpi_P^{t_f(q)})$. (If \cite{KO}'s coefficient ring $\cO\subseteq\widetilde\cO_P$ has ramification index $e$,
their exponent is $t_f(q)/e$; both sides of every identity in this chapter scale by $e$.)

(c) $f$ is $q$-new, so $\rho_f$ is ramified at $q$ (I2), $\rho_f(\tau)\ne1$ and $t_f(q)<\infty$.
\end{proof}
\noindent \emph{The ring $\widetilde\cO_P$.} $I_F/P$ is a one-dimensional complete Noetherian local domain ($P$ has height one), but need not be integrally closed: $I_F$ normal and $P$ of height one make the \emph{localisation} $(I_F)_P$ a DVR, not the quotient. We therefore pass to the normalisation $\widetilde\cO_P$ of $I_F/P$: it is finite over $I_F/P$ (complete Noetherian local rings are Nagata rings, Matsumura, \emph{Commutative ring theory}, \S32), local (Henselian) and normal of dimension one, hence a DVR, finite over $\cO$, and $I_F\to\widetilde\cO_P$ is a \emph{local} homomorphism (integral extension of the local ring $I_F/P$, so the maximal ideal contracts to the maximal ideal) which is in general not surjective. Lemma~\ref{lem:s06-S} base-changes along it, and $\cO\to\widetilde\cO_P$ is flat. \provisional{R1 (audit S06-M9): before, $\cO_P=I_F/P$ was claimed to be a DVR.}

\noindent (a) holds at every arithmetic point by local--global compatibility (Steinberg at $q$), and the arithmetic points are dense.
(b) Ramification indices of $\widetilde\cO_P$ over $\cO_{f,\lambda}$ rescale both sides equally.
